\section{Discussion}\label{sec:discussion}

\paragraph{Why not 2$\times$? Sharding split the rows, not the commits.}
Per post, the one database applied 46 rows in 1.02 commits. After sharding each database applied 21 or 24 rows
per post, about half, but still 1.03 commits per post: the author's shard commits \texttt{publish\_post()},
and the other shard commits \texttt{fanout()} for the followers who live there. With 40 random followers per
user, almost every post has followers on both shards. Commit-related waits (WALWrite, WalSync) were a large share
of the single database's saturation (Figure~\ref{fig:waits}), and the commit rate each Nano disk must sustain went from
587/s (one database at its peak) to 765/s per shard (two shards at their peak). Sharding therefore removed most of the
CPU and lock contention from inserting and indexing rows. It did not remove the per-commit cost. The
measured gain is 1.30--1.40$\times$, not the ideal 2$\times$ of an app whose writes stay with one owner (lecture
slide~9, e-commerce).

\paragraph{The cost of copy-at-write across shards.} The same design decision has a latency cost: one publish is
two sequential round trips (Section~\ref{sec:results}), so at low and medium load users see a p50 of about 155~ms instead of 78~ms.
The lecture accepts this ordering because the real record has to exist before its copies. Two fixes follow from the
lecture's own principles:
\begin{enumerate}[nosep,leftmargin=1.4em]
\item \emph{Combine many small writes into fewer} (slide~25): queue the follower copies and apply them per shard
      in batches, e.g.\ one \texttt{fanout()} transaction per shard every 50~ms carrying many posts. The post returns
      after one round trip, and the second shard's commit rate drops from one per post to one per batch.
\item \emph{Collect at read time} for authors with many followers (slide~50): no copies at all, so the publish is a
      one-shard, one-commit transaction, and the feed read asks both shards and merges.
\end{enumerate}
We did not implement these. They change the semantics (feeds become eventually consistent within the batch interval),
and the assignment asks for the lecture's ten steps.

\paragraph{Hot shard and hot owner.} With 400 users, hashing gave a 45/55 split, and the busier shard limited the
overloaded 256-client stage. More users would even this out; the lecture's 2{,}000-user test gave 1{,}008:992. One
very busy owner, the worksheet's Step~3 celebrity, cannot be split by hashing: all of its likes land on one shard.
Because likes are inserted as rows and never update a counter, they do not queue on a single row lock.

\paragraph{Revised capacity statement.} With two Nano shards, Campus Pulse sustained about 700 posts/s
($\approx$31{,}000--33{,}000 rows/s across both databases) with p99 $\le$ 400~ms, against 574 posts/s with p99 276~ms on one Nano database.
The M2 1M-user model needs about 17 posts/s (10\% of 167 write RPS). Both configurations cover that, so for
Campus Pulse at class scale sharding would be premature (``shard last'', slide~25). The numbers show what sharding
buys once a single database is saturated: about 1.3--1.4$\times$ for this fan-out-heavy write. Writes that touch one
owner (likes, sessions) should scale closer to 2$\times$, but our like test reached the API's $\approx$1{,}600~req/s ceiling
before either shard saturated (+21\%), so we did not measure their limit.

\paragraph{Confounds and limits.}
\begin{itemize}[nosep,leftmargin=1.2em]
\item \textbf{Shared free-tier hardware.} Nano projects run on shared hosts. Overloaded stages varied by up to 57\%
      between repeats (256 clients, two shards); unsaturated stages repeated within 1\% (128 clients, two shards:
      698 vs.\ 703 posts/s). We report repeats instead of averaging them.
\item \textbf{30~s measurement windows} (M2 used 60~s) because a one-database stage writes $\approx$13~KB per post and
      must stay below the 500~MB free-plan limit. Every stage starts from truncated tables.
\item \textbf{Host metrics have 60~s resolution and are not aligned with stages}, so we show them as a timeline and use
      the in-database 1~Hz samples for per-stage claims.
\item \textbf{Two different Postgres patch versions} (17.6 on shard\_0, 17.11 on shard\_1); same region, same compute size;
      round-trip times 73 and 75~ms.
\item \textbf{Load generator and API on one lab server}, 75~ms from Supabase. The API ceiling (about 1{,}600 req/s for
      requests that barely touch the database) is above every publish result we report. Hosted Render was not used
      for M6 because its single free instance tops out near 15 requests/s (M2).
\item \textbf{Closed-loop, uniform workload}: 40 followers per user and one fixed post body with three hashtags. Real
      follower counts are skewed, which would make the hot-owner effects larger.
\end{itemize}

\paragraph{Housekeeping.} The M4/M4b experiment schemas (about 415~MB) were exported to compressed CSV and dropped
from shard\_0 to make room; the M1/M2 \texttt{public} tables, the M5 \texttt{search} schema and the hosted app were
not modified, and M5's cron refresh was re-enabled after each series.
